% Zone „Das kennst du schon“ – aus bank/bedingte-wahrscheinlichkeit-und-bayes/zone.jsonl (zusammenbau v0.1)
\einheitenkopf[zone][Kennst du schon]{Das kennst du schon}
\setcounter{aufgabe}{0}

%% TODO Titel: Ich-kann-Satz fehlt in der Bank (Platzhalter: Kettenname) – Nr. 1
%% TODO Anweisung: ein Satz über der ersten Teilaufgabe fehlt in der Bank – Nr. 1
\begin{aufgabe}{Vierfeldertafel füllen und lesen (Rand, Feld, bedingte Angabe unterscheiden)}
\begin{teile}
\teil In einer Gruppe von 100 Personen bedeutet A: „trägt eine Brille“, B: „treibt Sport“. Ergänze die beiden fehlenden Felder der Vierfeldertafel.

\vierfeldertafel{A}{B}{30,10,40,,,60,45,55,100}
\leerfeld \quad \leerfeld
\teil Für die Ereignisse A und B gilt $P(A) = 0{,}3$ und $P(A \cap B) = 0{,}18$. Wie groß ist $P(A \cap \overline{B})$? \leerfeld
\end{teile}
\end{aufgabe}

%% TODO Titel: Ich-kann-Satz fehlt in der Bank (Platzhalter: Kettenname) – Nr. 2
%% TODO Anweisung: ein Satz über der ersten Teilaufgabe fehlt in der Bank – Nr. 2
\begin{aufgabe}{Baumdiagramm und Pfadregeln (Produkt entlang des Pfades, Summe über Pfade)}
\begin{teile}
\teil Aus einem Beutel wird zweimal gezogen; R heißt rot, B heißt blau. Wie groß ist die Wahrscheinlichkeit für erst rot, dann blau? Lies die Werte am Baum ab.

\baumzwei{R/0.4,B/0.6}{R/0.5,B/0.5}{R/0.3,B/0.7}
\leerfeld
\teil Erst wird eine von drei Karten gezogen: mit $\frac{1}{3}$ ein Ass, mit $\frac{2}{3}$ keins. Nach einem Ass gewinnt man mit $\frac{3}{4}$, sonst mit $\frac{1}{2}$. Wie groß ist die Wahrscheinlichkeit für Ass und Gewinn? \leerfeld
\end{teile}
\end{aufgabe}

%% TODO Titel: Ich-kann-Satz fehlt in der Bank (Platzhalter: Kettenname) – Nr. 3
%% TODO Anweisung: ein Satz über der ersten Teilaufgabe fehlt in der Bank – Nr. 3
\begin{aufgabe}{Anteile und Prozentsätze umrechnen, Quotienten von Anteilen bilden und kürzen}
\begin{teile}
\teil Schreibe $0{,}35$ als Prozentsatz. \leerfeld[\%]
\teil Berechne $\frac{0{,}14}{0{,}35}$ als Dezimalzahl. \leerfeld
\end{teile}
\end{aufgabe}

%% TODO Titel: Ich-kann-Satz fehlt in der Bank (Platzhalter: Kettenname) – Nr. 4
%% TODO Anweisung: ein Satz über der ersten Teilaufgabe fehlt in der Bank – Nr. 4
\begin{aufgabe}{Brüche mit Termen vergleichen (Zähler konstant, Nenner ändert sich) und einfache Bruchungleichungen lösen}
\begin{teile}
\teil Welcher Bruch ist größer: $\frac{3}{5}$ oder $\frac{3}{7}$? \leerfeld
\teil Wie ändert sich $\frac{4}{4 + x}$, wenn x von 1 auf 3 wächst? \leerfeld
\end{teile}
\end{aufgabe}

%% TODO Titel: Ich-kann-Satz fehlt in der Bank (Platzhalter: Kettenname) – Nr. 5
%% TODO Anweisung: ein Satz über der ersten Teilaufgabe fehlt in der Bank – Nr. 5
\begin{aufgabe}{Graphen lesen (Funktionswert zu Stelle, Stelle zu Funktionswert, Randwerte)}
\begin{teile}
\teil Lies am Graphen den Funktionswert $f(2)$ ab.

\begin{ksys}[xmin=0,xmax=6,ymin=0,ymax=5,ablesen] \funktion{0.5*\x+1}{f} \end{ksys}
$f(2) =$ \leerfeld
\teil Die Funktion h ist nur für $0 \le x \le 4$ gezeichnet. Lies ihre Werte am linken und am rechten Rand ab.

\begin{ksys}[xmin=0,xmax=4,ymin=0,ymax=5,ablesen] \funktionab{4/(\x+1)}{h}{0}{4} \end{ksys}
\leerfeld \quad \leerfeld
\end{teile}
\end{aufgabe}

%% TODO Titel: Ich-kann-Satz fehlt in der Bank (Platzhalter: Kettenname) – Nr. 6
%% TODO Anweisung: ein Satz über der ersten Teilaufgabe fehlt in der Bank – Nr. 6
\begin{aufgabe}{Vierfeldertafel füllen und lesen (Rand, Feld, bedingte Angabe unterscheiden) – Fehler finden}
\begin{teile}
\teil In einer Klasse bedeutet A: „trägt eine Brille“, B: „ist ein Mädchen“. Lea soll ablesen, wie viele Mädchen eine Brille tragen, und schreibt 15. Finde den Fehler.

\vierfeldertafel{A}{B}{9,6,15,7,8,15,16,14,30}
\end{teile}
\end{aufgabe}

%% TODO Titel: Ich-kann-Satz fehlt in der Bank (Platzhalter: Kettenname) – Nr. 7
%% TODO Anweisung: ein Satz über der ersten Teilaufgabe fehlt in der Bank – Nr. 7
\begin{aufgabe}{Vierfeldertafel füllen und lesen (Rand, Feld, bedingte Angabe unterscheiden) – selbst rechnen}
\begin{teile}
\teil In einer Klasse bedeutet A: „hat ein Haustier“, B: „ist ein Mädchen“. Wie viele Jungen haben kein Haustier? Lies in der Vierfeldertafel ab.

\vierfeldertafel{A}{B}{6,10,16,5,9,14,11,19,30}
\leerfeld
\end{teile}
\end{aufgabe}
